\documentclass[12pt,a4paper]{amsart}

\usepackage{amsmath,amssymb,amsthm}
\usepackage{mathtools}
\usepackage[utf8]{inputenc}
\usepackage[ngerman]{babel}
\usepackage{hyperref}
\usepackage{geometry}
\geometry{margin=2.5cm}
\usepackage[T1]{fontenc}

% -------------------------------------------------------
% Theorem-Umgebungen
% -------------------------------------------------------
\newtheorem{theorem}{Satz}[section]
\newtheorem{lemma}[theorem]{Lemma}
\newtheorem{corollary}[theorem]{Korollar}
\newtheorem{proposition}[theorem]{Proposition}
\theoremstyle{definition}
\newtheorem{definition}[theorem]{Definition}
\newtheorem{remark}[theorem]{Bemerkung}
\newtheorem{conjecture}[theorem]{Vermutung}

% -------------------------------------------------------
% Eigene Befehle
% -------------------------------------------------------
\newcommand{\N}{\mathbb{N}}
\newcommand{\Z}{\mathbb{Z}}
\newcommand{\Q}{\mathbb{Q}}
\newcommand{\R}{\mathbb{R}}
\newcommand{\C}{\mathbb{C}}
\newcommand{\Fp}{\mathbb{F}_p}
\DeclareMathOperator{\ord}{ord}
% Tate-Shafarevich-Gruppe: Sha als Mathoperator
\DeclareMathOperator{\Sha}{\mathfrak{Sh}}
\DeclareMathOperator{\rang}{rang}

% -------------------------------------------------------
\begin{document}
% -------------------------------------------------------

\title{Die Birch--Swinnerton-Dyer-Vermutung}

\author{Michael Fuhrmann}
\address{Specialist Mathematics Project, Build 84}
\email{specialist-maths@localhost}
\date{\today}

\begin{abstract}
Die Birch--Swinnerton-Dyer-Vermutung (BSD) ist eine der sieben
Millennium-Preisaufgaben des Clay Mathematics Institute.
Sie sagt eine tiefe Verbindung zwischen der Arithmetik einer elliptischen
Kurve $E/\Q$ (ihrem Rang $r$, Torsion, Periodenstruktur) und dem
analytischen Verhalten ihrer $L$-Funktion $L(E,s)$ bei $s=1$ vorher.

Wir formulieren die \emph{schwache BSD}: $L(E,1) = 0 \iff E(\Q)$ ist unendlich.
Wir formulieren die \emph{starke BSD}: $\ord_{s=1}L(E,s) = r$, und geben
die vollst\"andige Formel f\"ur den f\"uhrenden Koeffizienten mit $\Omega_E$,
den Tamagawa-Zahlen $c_p$, der Tate--Shafarevich-Gruppe $\Sha(E/\Q)$
und dem Regulator $R_E$.

Bekannte Resultate: R\"ange $0$ und $1$
(Coates--Wiles 1977, Kolyvagin 1990).
Numerische Evidenz: $\prod_{p \le X} \#E(\Fp)/p \sim C(\log X)^r$.
\end{abstract}

\maketitle
\tableofcontents

% ================================================================
\section{Einleitung}
% ================================================================

In den 1960er Jahren berechneten Birch und Swinnerton-Dyer \cite{BSD1965}
das Produkt $\prod_{p \le X} \#E(\Fp)/p$ f\"ur viele elliptische Kurven
und beobachteten numerisch ein Wachstum wie $C \cdot (\log X)^r$,
wobei $r$ der Rang von $E(\Q)$ ist.
Dies f\"uhrte zu der Vermutung, dass die Verschwindungsordnung von
$L(E,s)$ bei $s=1$ gleich dem Rang ist.

% ================================================================
\section{Schwache BSD: Rang und Verschwindungsordnung}
% ================================================================

\begin{conjecture}[Schwache BSD]
\label{conj:weak_bsd}
Sei $E$ eine elliptische Kurve \"uber $\Q$ mit Rang $r = \rang E(\Q)$.
Dann gilt
\[
  \ord_{s=1} L(E,s) \;=\; r.
\]
Insbesondere:
\begin{itemize}
  \item $L(E,1) \ne 0 \iff r = 0 \iff E(\Q)$ ist endlich.
  \item $L(E,1) = 0$ und $L'(E,1) \ne 0 \iff r = 1$.
\end{itemize}
\end{conjecture}

\begin{remark}[Aktueller Stand]
Die schwache BSD ist nur in Spezialf\"allen bewiesen:
\begin{itemize}
  \item Rang $0$: Coates--Wiles (1977) \cite{CoatesWiles1977} bewiesen $r=0$
    f\"ur $E$ mit komplexer Multiplikation und $L(E,1) \ne 0$.
  \item Kolyvagin (1990) \cite{Kolyvagin1990}: F\"ur modulare $E$ mit
    $L(E,1) \ne 0$ sind $E(\Q)$ und $\Sha(E/\Q)$ endlich.
    F\"ur $L'(E,1) \ne 0$ (und Vorzeichen $-1$) gilt $r = 1$.
  \item Gross--Zagier (1986): Unter bestimmten Bedingungen impliziert
    $L'(E,1) \ne 0$ die Existenz eines rationalen Punktes unendlicher Ordnung.
  \item Der allgemeine Fall ist offen.
\end{itemize}
\end{remark}

% ================================================================
\section{Die Tate--Shafarevich-Gruppe}
% ================================================================

\begin{definition}[Tate--Shafarevich-Gruppe]
\label{def:sha}
Sei $E/\Q$. Die \emph{Tate--Shafarevich-Gruppe} ist
\[
  \Sha(E/\Q) \;=\; \ker\!\Bigl(H^1(\Q, E) \to \prod_v H^1(\Q_v, E)\Bigr),
\]
wobei $v$ alle Stellen (Primzahlen $p$ und die archimedische Stelle $\infty$)
durchl\"auft.

Anschaulich: $\Sha(E/\Q)$ parametrisiert Torsoren unter $E$, die \"uber
jeder Vervollst\"andigung $\Q_v$ einen Punkt haben, aber keinen rationalen
Punkt \"uber $\Q$ -- das sind Versagen des lokalen-globalen Prinzips f\"ur $E$.
\end{definition}

\begin{proposition}[Grundeigenschaften von $\Sha$]
\label{prop:sha_basic}
\begin{enumerate}
  \item[\textup{(i)}] $\Sha(E/\Q)$ ist eine Torsions-abelsche Gruppe.
  \item[\textup{(ii)}] Die Cassels--Tate-Paarung ist alternierend
    (schiefsymmetrisch), also ist $|\Sha(E/\Q)|$ ein Quadrat (wenn endlich).
  \item[\textup{(iii)}] Unter BSD-Annahme ist $\Sha(E/\Q)$ endlich.
  \item[\textup{(iv)}] Kolyvagin bewies die Endlichkeit von $\Sha(E/\Q)$
    f\"ur analytischen Rang $\le 1$.
\end{enumerate}
\end{proposition}

% ================================================================
\section{Starke BSD: Die Formel f\"ur den f\"uhrenden Koeffizienten}
% ================================================================

\begin{conjecture}[Starke BSD]
\label{conj:strong_bsd}
Sei $E/\Q$ mit Rang $r$, reeller Periode $\Omega_E$, Regulator $R_E$,
endlicher Tate--Shafarevich-Gruppe $\Sha(E/\Q)$, Torsionsgruppe der
Ordnung $|E(\Q)_{\mathrm{tors}}|$ und Tamagawa-Zahlen
$c_p = |E(\Q_p)/E_0(\Q_p)|$ f\"ur Primzahlen $p$ schlechter Reduktion.
Dann gilt:
\begin{equation}
  \label{eq:strong_bsd}
  \lim_{s \to 1} \frac{L(E,s)}{(s-1)^r}
  \;=\;
  \frac{\Omega_E \cdot |\Sha(E/\Q)| \cdot R_E \cdot \prod_p c_p}
       {|E(\Q)_{\mathrm{tors}}|^2}.
\end{equation}
\end{conjecture}

\begin{remark}[Die Zutaten]
\begin{itemize}
  \item $\Omega_E = \int_{E(\R)} |\omega|$: die reelle Periode der
    N\'eron-Differentialform $\omega = dx/2y$.
  \item $R_E = \det(\langle P_i, P_j\rangle)$: der N\'eron--Tate-Regulator,
    Determinante der H\"ohenpaarungsmatrix
    f\"ur eine Basis $\{P_1,\ldots,P_r\}$ von $E(\Q)/E(\Q)_{\mathrm{tors}}$.
  \item $c_p$: die Tamagawa-Zahl bei $p$ (Anzahl der Komponenten des
    N\'eron-Modells bei $p$).
  \item $\Sha(E/\Q)$: Tate--Shafarevich-Gruppe (ihr Quadrat ist die Ordnung).
  \item $E(\Q)_{\mathrm{tors}}$: Torsionsuntergruppe (Paper~25).
\end{itemize}
\end{remark}

% ================================================================
\section{Numerische Evidenz: Das BSD-Produkt}
% ================================================================

\begin{definition}[BSD-Heuristikprodukt]
\label{def:bsd_product}
F\"ur $X > 0$ definieren wir das \emph{BSD-Heuristikprodukt}
\[
  B_E(X) \;=\; \prod_{\substack{p \le X\\ p\text{ gute Reduktion}}}
  \frac{\#E(\Fp)}{p}.
\]
\end{definition}

\begin{theorem}[Numerische Beobachtung von Birch--Swinnerton-Dyer]
\label{thm:bsd_numerical}
Birch und Swinnerton-Dyer \cite{BSD1965} berechneten $B_E(X)$ f\"ur viele
Kurven und beobachteten:
\[
  B_E(X) \;\sim\; C_E\,(\log X)^r \quad (X \to \infty),
\]
wobei $r = \rang E(\Q)$ und $C_E$ eine kurvenabh\"angige Konstante ist.
\end{theorem}

\begin{proof}[Heuristische Begr\"undung]
Mit Hasses Satz: $\#E(\Fp)/p = 1 + 1/p - a_p/p \approx 1 - a_p/p$.
Das Produkt $\prod_{p \le X}(1-a_p/p)$ h\"angt mit $L(E,1)$ via dem
Euler-Produkt zusammen.
Wenn $L(E,1) = 0$ (Rang $r \ge 1$), divergiert das Produkt logarithmisch
und ergibt das beobachtete $(\log X)^r$-Wachstum.
\end{proof}

% ================================================================
\section{Bekannte Resultate}
% ================================================================

\begin{theorem}[Coates--Wiles, 1977]
\label{thm:coates_wiles}
Sei $E/\Q$ eine elliptische Kurve mit komplexer Multiplikation durch
einen imagin\"ar-quadratischen K\"orper $K$.
Falls $L(E,1) \ne 0$, so ist $E(\Q)$ endlich (Rang $r = 0$).
\end{theorem}

\begin{proof}[Beweisskizze]
Coates und Wiles \cite{CoatesWiles1977} verwendeten eine $p$-adische
$L$-Funktion und die Arithmetik des CM-K\"orpers, um zu zeigen, dass
$L(E,1) \ne 0$ die Endlichkeit des $p^\infty$-Selmer-Raums impliziert.
\end{proof}

\begin{theorem}[Kolyvagin, 1990]
\label{thm:kolyvagin}
Sei $E/\Q$ eine modulare elliptische Kurve.
\begin{enumerate}
  \item[\textup{(i)}] Falls $L(E,1) \ne 0$, so $r = 0$ und
    $\Sha(E/\Q)$ ist endlich.
  \item[\textup{(ii)}] Falls $L'(E,1) \ne 0$, so $r = 1$ und
    $\Sha(E/\Q)$ ist endlich.
\end{enumerate}
\end{theorem}

\begin{proof}[Beweisskizze]
Kolyvagins Beweis \cite{Kolyvagin1990} verwendet \emph{Euler-Systeme}
von Heegner-Punkten: Punkte auf $E$ aus der CM-Theorie, die ein
normkompatibles System \"uber wachsenden K\"orpern bilden.
Schritte:
\begin{enumerate}
  \item Gross--Zagier (1986): Bei $L'(E,1) \ne 0$ ist der Heegner-Punkt
    auf $E$ ein Punkt unendlicher Ordnung.
  \item Kolyvaginscher Abstieg: Euler-System-Relationen liefern die
    vorhergesagte Struktur des Selmer-Raums.
  \item Endlichkeit von $\Sha$ folgt aus der Kolyvagin-System-Struktur.
\end{enumerate}
\end{proof}

% ================================================================
\section{Abstieg und die Selmer-Gruppe}
% ================================================================

\begin{remark}[Berechnung des Selmer-Raums und Abstieg]
Praktisch berechnet man den Rang durch einen \emph{2-Abstieg}:
Dies liefert den \emph{Selmer-Raum} $\mathrm{Sel}_2(E/\Q)$ in der exakten Folge
\[
  0 \to E(\Q)/2E(\Q) \to \mathrm{Sel}_2(E/\Q) \to \Sha(E/\Q)[2] \to 0.
\]
Die Abbildung $E(\Q)/2E(\Q) \to \mathrm{Sel}_2(E/\Q)$ ist die
\emph{Kummer-Einbettung}, die jeden Punkt $P$ durch die $2$-Isogenie-Bilder
lokal beschreibt.  Der Selmer-Raum ist ein endlicher $\mathbb{F}_2$-Vektorraum,
der lokal-global-l\"osbare Homomorphismen erfasst; seine Dimension gibt eine
obere Schranke f\"ur $r + \dim_{\mathbb{F}_2} \Sha(E/\Q)[2]$.

Die exakte algebraische Beziehung (ohne Sha$=0$-Annahme) lautet:
\[
  r + \dim_{\mathbb{F}_2} \Sha(E/\Q)[2]
  \;=\; \dim_{\mathbb{F}_2} \mathrm{Sel}_2(E/\Q)
        - \dim_{\mathbb{F}_2} E(\Q)_{\mathrm{tors}}[2].
\]
Dies folgt aus der exakten Folge $0 \to E(\Q)/2E(\Q) \to \mathrm{Sel}_2(E/\Q) \to \Sha(E/\Q)[2] \to 0$
und dem Struktursatz $E(\Q)/2E(\Q) \cong (\mathbb{Z}/2\mathbb{Z})^r \times E(\Q)_{\mathrm{tors}}[2]$
(da $E(\Q) \cong \mathbb{Z}^r \times E(\Q)_{\mathrm{tors}}$).
Unter der Annahme $\Sha(E/\Q) = 0$ vereinfacht sich dies zu
$r = \dim_{\mathbb{F}_2} \mathrm{Sel}_2(E/\Q) - \dim_{\mathbb{F}_2} E(\Q)_{\mathrm{tors}}[2]$.
Kolyvagin \cite{Kolyvagin1990} zeigte: wenn $L(E,1) \ne 0$ oder $L'(E,1) \ne 0$,
dann ist $\Sha(E/\Q)$ endlich und $r \le 1$.
\end{remark}

% ================================================================
\section{Schlussfolgerung}
% ================================================================

Die Birch--Swinnerton-Dyer-Vermutung sagt eine pr\"azise Verbindung zwischen
der Arithmetik von $E(\Q)$ (Rang, Regulator, Tate--Shafarevich-Gruppe)
und dem analytischen Verhalten von $L(E,s)$ bei $s=1$ vorher.

Die schwache Form (Rang = analytischer Rang) ist nur f\"ur analytischen
Rang $\le 1$ bewiesen (Kolyvagin 1990).
Die starke Form~\eqref{eq:strong_bsd} wurde numerisch f\"ur tausende Kurven
verifiziert, ist aber im Allgemeinen unbewiesen.

BSD ist eine der sieben Millennium-Preisaufgaben; ihre L\"osung w\"urde
ein fundamentales Verst\"andnis rationaler Punkte auf algebraischen
Variet\"aten erm\"oglichen.

% ================================================================
\begin{thebibliography}{10}

\bibitem{BSD1965}
B.~J.~Birch und H.~P.~F.~Swinnerton-Dyer,
Notes on elliptic curves, II,
\emph{J.\ Reine Angew.\ Math.}\ \textbf{218} (1965), 79--108.

\bibitem{CoatesWiles1977}
J.~Coates und A.~Wiles,
On the conjecture of Birch and Swinnerton-Dyer,
\emph{Invent.\ Math.}\ \textbf{39} (1977), 223--251.

\bibitem{Kolyvagin1990}
V.~A.~Kolyvagin,
Euler systems,
in: \emph{The Grothendieck Festschrift}, Vol.~II,
Progr.\ Math.~87, Birkh\"auser, 1990, S.~435--483.

\bibitem{GrossZagier1986}
B.~H.~Gross und D.~B.~Zagier,
Heegner points and derivatives of $L$-series,
\emph{Invent.\ Math.}\ \textbf{84} (1986), 225--320.

\bibitem{Silverman1986}
J.~H.~Silverman,
\emph{The Arithmetic of Elliptic Curves},
Graduate Texts in Mathematics~106,
Springer, New York, 1986.

\bibitem{Wiles1995}
A.~Wiles,
Modular elliptic curves and Fermat's last theorem,
\emph{Ann.\ of Math.}\ (2) \textbf{142} (1995), 443--551.

\bibitem{Milne2006}
J.~S.~Milne,
\emph{Elliptic Curves},
BookSurge Publishing, 2006.
Verf\"ugbar unter: \texttt{https://www.jmilne.org/math/Books/EC.pdf}.

\end{thebibliography}

\end{document}
